\documentclass[11pt]{article}
\usepackage{mathblatt}


\blattfuss{Kurvenuntersuchung}{Lösungen}
\begin{document}
\noindent{\large\bfseries Kurvenuntersuchung \textperiodcentered{} Lösungen}\par\medskip
\begin{pfloesung}{}
\lz{1.}{$5x^4$}{}{}
\end{pfloesung}
\begin{pfloesung}{}
\lz{2.}{$4$}{}{}
\end{pfloesung}
\begin{pfloesung}{}
\lz{3.}{$H(1 | 2)$}{}{}
\end{pfloesung}
\begin{pfloesung}{}
\lz{4.}{$7$}{}{}
\end{pfloesung}
\begin{pfloesung}{}
\lz{5.}{nach unten geöffnet}{}{}
\end{pfloesung}
\begin{pfloesung}{}
\lz{6.}{der Graph steigt in $A$}{}{}
\end{pfloesung}
\begin{pfloesung}{}
\lz{7.}{$x_1 = -7$, $x_2 = 7$}{}{}
\end{pfloesung}
\begin{pfloesung}{}
\lz{8.}{Faktoren null setzen}{}{}
\end{pfloesung}
\begin{pfloesung}{}
\lz{9.}{Stelle gegeben, Wert gesucht}{}{}
\end{pfloesung}
\begin{pfloesung}{}
\lz{10.}{$x_1 = 0$, $x_2 = -3$, $x_3 = 3$}{$x(x^2 - 9) = 0$}{}
\end{pfloesung}
\begin{pfloesung}{}
\lz{11.}{$x_1 = 1$, $x_2 = 5$}{}{}
\end{pfloesung}
\begin{pfloesung}{}
\lz{12.}{$x_1 = -3$, $x_2 = -2$, $x_3 = 2$, $x_4 = 3$}{$z = 4$ oder $z = 9$}{}
\end{pfloesung}
\begin{pfloesung}{}
\lz{13.}{$x_1 = -2$, $x_2 = 2$}{$e^x > 0 \Rightarrow x^2 - 4 = 0$}{}
\end{pfloesung}
\begin{pfloesung}{}
\lz{14.}{$N_1(-2 | 0)$, $N_2(0 | 0)$, $N_3(3 | 0)$}{}{}
\end{pfloesung}
\begin{pfloesung}{{\small\color{mbgrau}Abi GK 2025 · A1.4 a}}
\lz{15.}{f(1) = 2e $-$ 2e = 0}{}{}
\end{pfloesung}
\begin{pfloesung}{{\small\color{mbgrau}Abi GK 2026 · B2.2 g}}
\lz{16.}{80 °C}{$k(0) = 60 + 20$}{}
\end{pfloesung}
\begin{pfloesung}{}
\lz{17.}{$f(2) = -2$, $P(2 | -2)$}{$f(2) = 2^3 - 4 \cdot 2^2 + 6$; $f(2) = 8 - 16 + 6 = -2$}{}
\end{pfloesung}
\begin{pfloesung}{{\small\color{mbgrau}Abi GK 2022 · A1.1 a}}
\lz{18.}{x = 0, x = 2, x = 6}{$f(x) = x(x - 2)(x - 6)$}{}
\end{pfloesung}
\begin{pfloesung}{{\small\color{mbgrau}Abi GK 2026 · B2.2 a}}
\lz{19.}{e$^{x}$ > 0 für alle x $\Rightarrow$ 3e$^{x}$ + 1 > 1 > 0 $\Rightarrow$ keine Nullstelle; Wertemenge ]1; $\infty$[}{}{}
\end{pfloesung}
\begin{pfloesung}{}
\lz{20.}{$x_1 = 3$, $x_2 = -5$}{$(x - 3)(x + 5) = 0$; $x - 3 = 0 \Rightarrow x_1 = 3$}{}
\end{pfloesung}
\begin{pfloesung}{{\small\color{mbgrau}Abi GK 2023 · B2.2 a}\par\noindent{\small $f'(x) = 3x^{2} - 24x + 45$}\par\noindent{\small $f''(x) = 6x - 24$}\par\noindent{\small $g'(x) = 1{,}2(x - 5)^{2}$}}
\lz{21.}{(0; $-$50); genau zwei Nullstellen x = 2 und x = 5 (doppelt)}{$f(x) = (x - 5)^{2} \cdot  (x - 2)$}{}
\end{pfloesung}
\begin{pfloesung}{{\small\color{mbgrau}Abi GK 2024 · B2.2 b}\par\noindent{\small $f''(x) = 6x^{2} - 8$}}
\lz{22.}{($-$1; 0) wegen Achsensymmetrie; weitere Schnittpunkte ($-\sqrt{7}$; 0), ($\sqrt{7}$; 0)}{mit a); u = x²: 0,5u² $-$ 4u + 3,5 = 0 $\Rightarrow$ u = 1 oder u = 7}{}
\end{pfloesung}
\begin{pfloesung}{{\small\color{mbgrau}Abi GK 2023 · B2.1 b}\par\noindent{\small $f''(x) = (0{,}5x^{2} + 2x - 1) \cdot  e^x$}}
\lz{23.}{(0; $-$2), ($-$2; 0), (2; 0); R(0; 2)}{$f(0) = -2$; $x^{2} - 4 = 0$}{}
\end{pfloesung}
\begin{pfloesung}{{\small\color{mbgrau}Abi GK 2024 · B2.1 a}}
\lz{24.}{Nullstelle x = 2; (0; $-$2e³) $\approx$ (0; $-$40,2)}{$f(0) = -2e^{3} \approx  -40{,}17$}{}
\end{pfloesung}
\begin{pfloesung}{{\small\color{mbgrau}Abi GK 2022 · B2.1 a}\par\noindent{\small $h'(t) = \tfrac{27}{40} t^{2} - 27t + \tfrac{405}{2}$}\par\noindent{\small $h''(t) = \tfrac{27}{20} t - 27$}}
\lz{25.}{($-$2; 0), (0; 2); f'(x) = 0 $\Rightarrow$ x = $-$1, Vorzeichenwechsel von + nach $-$ $\Rightarrow$ Hochpunkt ($-$1; e)}{$f(0) = 2$; $f(-1) = e \approx  2{,}72$}{}
\end{pfloesung}
\begin{pfloesung}{}
\lz{26.}{nachweisen: $x = 2$ einsetzen}{}{}
\end{pfloesung}
\begin{pfloesung}{}
\lz{27.}{gefragt ist nur, wo die Tangente waagerecht ist}{}{}
\end{pfloesung}
\begin{pfloesung}{}
\lz{28.}{$-18 < 0$, also Hochpunkt}{$f'(x) = 3x^2 - 27$, $f'(-3) = 27 - 27 = 0$; $f''(x) = 6x$, $f''(-3)$}{}
\end{pfloesung}
\begin{pfloesung}{}
\lz{29.}{$0$}{$f'(x) = 3x^2 - 12$; $f'(2) = 3 \cdot 4 - 12$}{}
\end{pfloesung}
\begin{pfloesung}{}
\lz{30.}{$-12 < 0$: Hochpunkt $H(4 | 27)$}{$f'(x) = -3x^2 + 12x = -3x(x - 4) = 0$ für $x_1 = 0$, $x_2 = 4$}{}
\end{pfloesung}
\begin{pfloesung}{}
\lz{31.}{$32 > 0$: Tiefpunkte $T_1(-2 | -16)$ und $T_2(2 | -16)$}{$f'(x) = 4x^3 - 16x = 4x(x^2 - 4) = 0$ für $x_1 = 0$, $x_2 = -2$, $x_3 = 2$}{}
\end{pfloesung}
\begin{pfloesung}{{\small\color{mbgrau}Abi GK 2023 · A1.3 a}}
\lz{32.}{x = 0 und x = 2}{$f'(x) = 3x^{2} - 6x$}{}
\end{pfloesung}
\begin{pfloesung}{}
\lz{33.}{$x_1 = 3$; $x_2 = 1$}{$f'(x) = x^2 - 4x + 3$; $x^2 - 4x + 3 = 0$}{}
\end{pfloesung}
\begin{pfloesung}{}
\lz{34.}{$5$, also ist $H(2 | 5)$ Hochpunkt}{$f'(x) = -3x^2 + 6x$, $f'(2) = 0$; $f''(x) = -6x + 6$, $f''(2) = -6 < 0$}{}
\end{pfloesung}
\begin{pfloesung}{{\small\color{mbgrau}Abi GK 2025 · B2.2 b}\par\noindent{\small $f'(x) = (1 - x) \cdot  e^x$}}
\lz{35.}{Hochpunkt (1; e)}{$f'(x) = 0 \Rightarrow  x = 1$}{}
\end{pfloesung}
\begin{pfloesung}{{\small\color{mbgrau}Abi GK 2022 · B2.2 b}\par\noindent{\small $f'(x) = -\tfrac{1}{2} x^{2} + x$}\par\noindent{\small $f''(x) = -x + 1$}}
\lz{36.}{Tiefpunkt (0; 0), Hochpunkt (2; $\tfrac{2}{3}$)}{}{}
\end{pfloesung}
\begin{pfloesung}{{\small\color{mbgrau}Abi GK 2025 · B2.1 b}\par\noindent{\small $f''(x) = \tfrac{12}{25} x$}}
\lz{37.}{Hochpunkt ($-$2,5; 2,5), Tiefpunkt (2,5; $-$2,5); Gerade durch beide: y = $-$x $\Rightarrow$ Winkelhalbierende des II. und IV. Quadranten}{$f'(x) = \tfrac{6}{25} x^{2} - \tfrac{3}{2} = 0 \Rightarrow  x = \pm 2{,}5$}{}
\end{pfloesung}
\begin{pfloesung}{{\small\color{mbgrau}Abi GK 2023 · B2.2 b}}
\lz{38.}{H(3; 4), T(5; 0)}{Kopf bei Nr. 21.}{}
\end{pfloesung}
\begin{pfloesung}{}
\lz{39.}{$H(2 | \mathrm{e}^{2})$, also $H(2 | 7{,}39)$}{$f'(x) = -\mathrm{e}^{x} + (3 - x)\,\mathrm{e}^{x} = (2 - x)\,\mathrm{e}^{x} = 0$; da $\mathrm{e}^{x} > 0$, ist $x = 2$}{}
\end{pfloesung}
\begin{pfloesung}{}
\lz{40.}{$y = -x$}{$f'(x) = \tfrac{3}{8}x^2 - 1{,}5 = 0 \Rightarrow x^2 = 4$; $f''(x) = \tfrac{3}{4}x$: $H(-2 | 2)$, $T(2 | -2)$}{}
\end{pfloesung}
\begin{pfloesung}{{\small\color{mbgrau}Abi GK 2022 · B2.1 j}}
\lz{41.}{h'(10) = 0, h''(10) = $-$13,5 < 0 $\Rightarrow$ Maximum; h(10) = 900 m}{Kopf bei Nr. 25.}{}
\end{pfloesung}
\begin{pfloesung}{{\small\color{mbgrau}Abi GK 2024 · B2.2 c}}
\lz{42.}{Hochpunkt (0; 3,5); Tiefpunkte ($-$2; $-$4,5), (2; $-$4,5)}{Kopf bei Nr. 22.; $f'(x) = 2x^{3} - 8x = 2x(x - 2)(x + 2)$}{}
\end{pfloesung}
\begin{pfloesung}{{\small\color{mbgrau}Abi GK 2023 · B2.1 c}}
\lz{43.}{f'($-$1 + $\sqrt{5}$) = 0 $\Rightarrow$ waagerechte Tangente; f($-$1 + $\sqrt{5}$) = (1 $-$ $\sqrt{5}$) · e$^{√5 - 1}$ $\approx$ $-$4,25}{Kopf bei Nr. 23.; $(-1 + \sqrt{5})^{2} = 6 - 2\sqrt{5}$}{}
\end{pfloesung}
\begin{pfloesung}{{\small\color{mbgrau}Abi GK 2024 · B2.1 b}}
\lz{44.}{Hochpunkt (4; 2e) $\approx$ (4; 5,44)}{$f'(x) = 0 \Rightarrow  x = 4$; Vorzeichenwechsel von + nach $-$}{}
\end{pfloesung}
\begin{pfloesung}{}
\lz{45.}{der Wert $f'''(1) = 6$ ist gegeben}{}{}
\end{pfloesung}
\begin{pfloesung}{}
\lz{46.}{gefragt ist die größte Wachstumsrate}{}{}
\end{pfloesung}
\begin{pfloesung}{}
\lz{47.}{Die Temperatur fällt von $90$ °C zunächst schnell, dann immer langsamer und nähert sich $20$ °C, der Raumtemperatur}{}{}
\end{pfloesung}
\begin{pfloesung}{}
\lz{48.}{$W(1 | 3)$}{$f''(x) = 6x - 6 = 0$ für $x = 1$; $f'''(1) = 6 \ne 0$}{}
\end{pfloesung}
\begin{pfloesung}{}
\lz{49.}{$W_1(-1 | -2)$, $W_2(1 | -2)$}{$f''(x) = 12x^2 - 12 = 0 \Rightarrow x^2 = 1 \Rightarrow x_1 = -1$, $x_2 = 1$}{}
\end{pfloesung}
\begin{pfloesung}{}
\lz{50.}{$6$}{$f'(x) = 3x^2 + 6x - 5$; $f''(x) = 6x + 6$}{}
\end{pfloesung}
\begin{pfloesung}{}
\lz{51.}{$-3$, also ist $W(2 | -3)$ Wendepunkt}{$f''(x) = 6x - 12$, $f''(2) = 0$; $f'''(2) = 6 \ne 0$}{}
\end{pfloesung}
\begin{pfloesung}{{\small\color{mbgrau}Abi GK 2023 · B2.2 c}}
\lz{52.}{W(4; 2)}{Kopf bei Nr. 21.}{}
\end{pfloesung}
\begin{pfloesung}{{\small\color{mbgrau}Abi GK 2024 · B2.1 c}}
\lz{53.}{f''(x) = (x/4 $-$ $\tfrac{3}{2}$) · e$^{-x/2 + 3}$ mit Produktregel; f''(6) = 0 mit Vorzeichenwechsel, f(6) = 4 $\Rightarrow$ Wendepunkt (6; 4)}{$f'''(6) = \tfrac{1}{4}$}{}
\end{pfloesung}
\begin{pfloesung}{{\small\color{mbgrau}Abi GK 2026 · B2.1 c}}
\lz{54.}{W1($-$2; $\tfrac{2}{3}$), W2(0; 2)}{$f''(x) = x^{2} + 2x = x(x + 2)$; $f'''(x) = 2x + 2: f'''(0) = 2, f'''(-2) = -2$}{}
\end{pfloesung}
\begin{pfloesung}{}
\lz{55.}{$W(2 | 4\,\mathrm{e}^{-1})$, also $W(2 | 1{,}47)$}{$f'(x) = -0{,}5x\,\mathrm{e}^{-0{,}5x}$}{}
\end{pfloesung}
\begin{pfloesung}{}
\lz{56.}{der Wendepunkt ist kein Richtungswechsel}{Bei $t = 10$ steigt der Ballon am schnellsten}{}
\end{pfloesung}
\begin{pfloesung}{{\small\color{mbgrau}Abi GK 2022 · B2.1 l}}
\lz{57.}{Wendepunkt (t = 20): Zeitpunkt der größten Geschwindigkeit beim Aufstieg}{Kopf bei Nr. 25.; $h''(t) = 0 \Rightarrow  t = 20$}{}
\end{pfloesung}
\begin{pfloesung}{{\small\color{mbgrau}Abi GK 2022 · A1.2 a}}
\lz{58.}{a = 0 (Tiefpunkt bei 0,2); b $\approx$ 0,5 (jede Wendestelle, auch 1,2 oder 1,8)}{}{}
\end{pfloesung}
\begin{pfloesung}{{\small\color{mbgrau}Abi GK 2022 · A1.2 b}}
\lz{59.}{P $\approx$ (1,8; 0,22); Wendepunkt (Krümmungswechsel) im fallenden Abschnitt zwischen Hochpunkt ($\approx$ 1,6) und Tiefpunkt ($\approx$ 2) $\Rightarrow$ f''(x\_P) = 0 und f'(x\_P) < 0}{mit a)}{}
\end{pfloesung}
\begin{pfloesung}{{\small\color{mbgrau}Abi GK 2023 · B2.1 f}}
\lz{60.}{W1($-$4,45; 0,09), W2(0,45; $-$2,98)}{Kopf bei Nr. 23.; $x = -2 \pm  \sqrt{6}$}{}
\end{pfloesung}
\begin{pfloesung}{}
\lz{61.}{die Tangente in $P$ geht nach rechts oben}{}{}
\end{pfloesung}
\begin{pfloesung}{}
\lz{62.}{$A$: plus, $B$: null, $C$: minus}{}{}
\end{pfloesung}
\begin{pfloesung}{}
\lz{63.}{für $4 \le x \le 6$ liegt $f'(x) = 2x - 6$ zwischen $2$ und $6$, also positiv}{}{}
\end{pfloesung}
\begin{pfloesung}{}
\lz{64.}{$x_1 = -1$ und $x_2 = 3$}{$f'(x) = 3x^2 - 6x - 9 = 0$ für}{}
\end{pfloesung}
\begin{pfloesung}{}
\lz{65.}{$f$ steigt auf $]-\infty; -2]$, fällt auf $[-2; 1]$ und steigt auf $[1; \infty[$}{}{}
\end{pfloesung}
\begin{pfloesung}{}
\lz{66.}{Der Hochpunkt $H(2 | 3)$: waagerechte Tangente und rechtsgekrümmt}{$f' = 0$; $f'' < 0$}{}
\end{pfloesung}
\begin{pfloesung}{}
\lz{67.}{$-15 < 0$: $f$ fällt auf $[3; \infty[$}{$f'(x) = -3x^2 + 6x + 9 = 0$ für $x_1 = -1$ und $x_2 = 3$}{}
\end{pfloesung}
\begin{pfloesung}{{\small\color{mbgrau}Abi GK 2026 · A1.7 a}}
\lz{68.}{f' wechselt an seiner einzigen Nullstelle $-$3 das Vorzeichen von $-$ nach + $\Rightarrow$ f fällt für x $\le$ $-$3 und steigt für x $\ge$ $-$3}{}{}
\end{pfloesung}
\begin{pfloesung}{{\small\color{mbgrau}Abi GK 2022 · B2.1 c}}
\lz{69.}{Graph II ist f'; Nullstelle bei x = $-$1 (Hochstelle von f), negativ für x > $-$1, weil f dort fällt}{Kopf bei Nr. 25.; mit a)}{}
\end{pfloesung}
\begin{pfloesung}{{\small\color{mbgrau}Abi GK 2026 · B2.1 h}}
\lz{70.}{Q(2; $-\tfrac{4}{3}$): tiefster Punkt des unteren Schienenverlaufs rechts, am Übergang in den rechten Halbkreis}{$g'(x) = x^{2} + 2x = 0 \Rightarrow  x = -2 (x \neq  0) \Rightarrow  x_Q = 2$; $y_Q = -g(-2) = -\tfrac{4}{3}$}{}
\end{pfloesung}
\begin{pfloesung}{}
\lz{71.}{links vom Tiefpunkt fällt der Graph, also ist $f'(1{,}9) < 0$}{}{}
\end{pfloesung}
\begin{pfloesung}{}
\lz{72.}{$f' > 0$ für $x < -1$, $f' < 0$ zwischen $-1$ und $3$, $f' > 0$ für $x > 3$}{Nullstellen von $f'$: $x = -1$ und $x = 3$}{}
\end{pfloesung}
\begin{pfloesung}{{\small\color{mbgrau}Abi GK 2024 · A1.7 a}\par\noindent{\small $f'(x) = 0{,}5 \cdot  e^{0{,}5x}$}}
\lz{73.}{f'(x) = 0,5 · e$^{0{,}5x}$ > 0 für alle x $\Rightarrow$ streng monoton steigend}{}{}
\end{pfloesung}
\begin{pfloesung}{{\small\color{mbgrau}Abi GK 2023 · B2.2 f}}
\lz{74.}{genau eine Nullstelle von g' (x = 5) $\Rightarrow$ genau eine waagerechte Tangente; g' $\ge$ 0 ohne Vorzeichenwechsel $\Rightarrow$ kein Extrempunkt (Sattelpunkt)}{Kopf bei Nr. 21.}{}
\end{pfloesung}
\begin{pfloesung}{{\small\color{mbgrau}Abi GK 2023 · B2.1 e}}
\lz{75.}{streng monoton fallend auf [0; $-$1 + $\sqrt{5}$], streng monoton steigend auf [$-$1 + $\sqrt{5}$; $\infty$)}{Kopf bei Nr. 23.; mit c); Tiefstelle $-$1 + $\sqrt{5}$ $\approx$ 1,24}{}
\end{pfloesung}
\begin{pfloesung}{{\small\color{mbgrau}Abi GK 2023 · B2.1 d}}
\lz{76.}{I falsch; Tiefpunkt hat den kleinsten Wert seiner Umgebung $\Rightarrow$ f($-$0,9 + $\sqrt{5}$) > f($-$1 + $\sqrt{5}$). II wahr; rechts der Tiefstelle steigt G $\Rightarrow$ f'($-$0,9 + $\sqrt{5}$) > 0}{Kopf bei Nr. 23.; mit c); $f(-0{,}9 + \sqrt{5}) \approx  -4{,}21$; $f'(-0{,}9 + \sqrt{5}) \approx  0{,}87$}{}
\end{pfloesung}
\begin{pfloesung}{}
\lz{77.}{Grad ungerade, Koeffizient negativ}{}{}
\end{pfloesung}
\begin{pfloesung}{}
\lz{78.}{der Faktor $e^{-3x}$}{}{}
\end{pfloesung}
\begin{pfloesung}{}
\lz{79.}{für $x \to -\infty$: $f(x) \to +\infty$}{}{}
\end{pfloesung}
\begin{pfloesung}{}
\lz{80.}{für $x \to -\infty$: $f(x) \to 0$}{}{}
\end{pfloesung}
\begin{pfloesung}{}
\lz{81.}{für $x \to -\infty$: $f(x) \to +\infty$}{}{}
\end{pfloesung}
\begin{pfloesung}{}
\lz{82.}{für $x \to -\infty$: $f'(x) \to +\infty$}{}{}
\end{pfloesung}
\begin{pfloesung}{}
\lz{83.}{$f(x) \to 0$: der e-Faktor fällt schneller, als $x + 4$ wächst}{}{}
\end{pfloesung}
\begin{pfloesung}{}
\lz{84.}{für $x \to +\infty$: $f(x) \to -\infty$}{}{}
\end{pfloesung}
\begin{pfloesung}{{\small\color{mbgrau}Abi GK 2022 · B2.1 b}}
\lz{85.}{Grenzwert 0; Graph nähert sich der x-Achse von oben}{Kopf bei Nr. 25.}{}
\end{pfloesung}
\begin{pfloesung}{{\small\color{mbgrau}Abi GK 2022 · B2.2 a}}
\lz{86.}{x $\to$ +$\infty$: f(x) $\to$ $-\infty$; x $\to$ $-\infty$: f(x) $\to$ +$\infty$}{Kopf bei Nr. 36.; Leitterm $-\tfrac{1}{6}$ x³}{}
\end{pfloesung}
\begin{pfloesung}{{\small\color{mbgrau}Abi GK 2026 · B2.1 a}}
\lz{87.}{x $\to$ $-\infty$: f'(x) $\to$ $-\infty$; x $\to$ +$\infty$: f'(x) $\to$ +$\infty$}{Leitterm $\tfrac{1}{3}$ x³}{}
\end{pfloesung}
\begin{pfloesung}{{\small\color{mbgrau}Abi GK 2025 · B2.2 a}}
\lz{88.}{Nullstelle 2; x $\to$ $-\infty$: f(x) $\to$ 0; x $\to$ +$\infty$: f(x) $\to$ $-\infty$}{Kopf bei Nr. 35.}{}
\end{pfloesung}
\begin{pfloesung}{}
\lz{89.}{für $x \to -\infty$: $f(x) \to 0$}{}{}
\end{pfloesung}
\begin{pfloesung}{{\small\color{mbgrau}Abi GK 2023 · B2.1 a}}
\lz{90.}{x $\to$ +$\infty$: f(x) $\to$ +$\infty$; x $\to$ $-\infty$: f(x) $\to$ 0}{Kopf bei Nr. 23.}{}
\end{pfloesung}
\begin{pfloesung}{}
\lz{91.}{achsensymmetrisch}{}{}
\end{pfloesung}
\begin{pfloesung}{}
\lz{92.}{punktsymmetrisch}{}{}
\end{pfloesung}
\begin{pfloesung}{}
\lz{93.}{prüfen: nur gerade Exponenten}{}{}
\end{pfloesung}
\begin{pfloesung}{}
\lz{94.}{$f(x)$: achsensymmetrisch zur y-Achse}{$f(-x) = (-x)^2 \cdot e^{-(-x)^2} = x^2 \cdot e^{-x^2}$}{}
\end{pfloesung}
\begin{pfloesung}{}
\lz{95.}{$S_y(0 | 7)$}{}{}
\end{pfloesung}
\begin{pfloesung}{}
\lz{96.}{bei Punktsymmetrie wird aus dem Tiefpunkt ein Hochpunkt}{}{}
\end{pfloesung}
\begin{pfloesung}{{\small\color{mbgrau}Abi GK 2024 · A1.1 a}}
\lz{97.}{nur ungerade Exponenten $\Rightarrow$ punktsymmetrisch zum Ursprung}{}{}
\end{pfloesung}
\begin{pfloesung}{{\small\color{mbgrau}Abi GK 2024 · B2.2 a}}
\lz{98.}{nur gerade Exponenten (f($-$x) = f(x)) $\Rightarrow$ achsensymmetrisch zur y-Achse}{Kopf bei Nr. 22.}{}
\end{pfloesung}
\begin{pfloesung}{{\small\color{mbgrau}Abi GK 2025 · B2.1 a}}
\lz{99.}{f($-$x) = $-\tfrac{2}{25}$ x³ + $\tfrac{3}{2}$ x = $-$f(x) $\Rightarrow$ punktsymmetrisch; Nullstellen 0, $\pm$5$\sqrt{3}$/2 $\approx$ $\pm$4,33}{Kopf bei Nr. 37.; $f(x) = x \cdot  (\tfrac{2}{25} x^{2} - \tfrac{3}{2})$; $x^{2} = \tfrac{75}{4}$}{}
\end{pfloesung}
\begin{pfloesung}{}
\lz{100.}{innen: x-Richtung}{}{}
\end{pfloesung}
\begin{pfloesung}{}
\lz{101.}{$4$, das Plus heißt nach links}{$2x + 8 = 2 \cdot (x + 4)$}{}
\end{pfloesung}
\begin{pfloesung}{}
\lz{102.}{$2$, Verschiebung um $2$ nach oben}{$a = 3$, Streckung in y-Richtung mit dem Faktor $3$}{}
\end{pfloesung}
\begin{pfloesung}{}
\lz{103.}{Verschiebung um $2$ nach links}{}{}
\end{pfloesung}
\begin{pfloesung}{}
\lz{104.}{$12$}{$f'(2) = g'(2)$}{}
\end{pfloesung}
\begin{pfloesung}{}
\lz{105.}{$H'(1 | 21)$}{$g(x) = (x - 3)^3 - 12(x - 3) + 5$}{}
\end{pfloesung}
\begin{pfloesung}{}
\lz{106.}{$4(x - 1)^2 - 3$: Streckfaktor $4$, Verschiebung um $1$ nach rechts und um $3$ nach unten}{}{}
\end{pfloesung}
\begin{pfloesung}{}
\lz{107.}{$P'(-2 | 1)$}{$y = -0{,}5x$}{}
\end{pfloesung}
\begin{pfloesung}{{\small\color{mbgrau}Abi GK 2026 · B2.2 f}}
\lz{108.}{Spiegelung an der y-Achse und Streckung in x-Richtung mit Faktor 400}{}{}
\end{pfloesung}
\begin{pfloesung}{{\small\color{mbgrau}Abi GK 2024 · A1.4 b}}
\lz{109.}{falsch; Verschiebung in y-Richtung erhält Steigungen, aber g'(0) = 2 $\neq$ h'(0) = 1}{mit a); $h'(0) = 1$}{}
\end{pfloesung}
\begin{pfloesung}{{\small\color{mbgrau}Abi GK 2023 · B2.2 j}}
\lz{110.}{c = 4, d = 55; Graph von f um 55 nach oben verschoben, dann in y-Richtung mit Faktor 4 gestreckt}{Kopf bei Nr. 21.; $f(x) + 55 = x^{3} - 12x^{2} + 45x + 5$}{}
\end{pfloesung}
\begin{pfloesung}{{\small\color{mbgrau}Abi GK 2024 · A1.7 b}}
\lz{111.}{h(x) = f(x + 2) = e$^{0{,}5x + 1}$ $-$ e}{Kopf bei Nr. 73.; Nullstelle von f: x = 2}{}
\end{pfloesung}
\begin{pfloesung}{{\small\color{mbgrau}Abi GK 2023 · B2.1 i}}
\lz{112.}{z. B. g(x) = f($-$x) = 0,5 · (x² $-$ 4) · e$^{-x}$ (IV. Quadrant) und h(x) = $-$f(x) = $-$0,5 · (x² $-$ 4) · e$^{x}$ (II. Quadrant); k(x) = $-$f($-$x) (I. Quadrant)}{Kopf bei Nr. 23.}{}
\end{pfloesung}
\begin{pfloesung}{{\small\color{mbgrau}Abi GK 2025 · B2.1 c}}
\lz{113.}{g* $\perp$ g durch den Ursprung $\Rightarrow$ Steigung $-$1/($\tfrac{1}{2}$) = $-$2 $\Rightarrow$ g*(x) = $-$2x; A*($-$2,5; 5)}{Kopf bei Nr. 37.; Drehung um 90°: (x | y) $\to$ ($-$y | x)}{}
\end{pfloesung}
\begin{pfloesung}{{\small\color{mbgrau}Abi GK 2026 · A1.4 b}}
\lz{114.}{nein; Graph von g entsteht aus Graph von f durch Streckung in x-Richtung (Faktor 3) und y-Richtung (Faktor 2), beide > 1 $\Rightarrow$ Strecke länger}{mit a); $\mid AB\mid  = \sqrt{\pi ^{2} + 4} \approx  3{,}72$; bei g: $\sqrt{9 \pi ² + 16}$ $\approx$ 10,24}{}
\end{pfloesung}
\begin{pfloesung}{{\small\color{mbgrau}Abi GK 2026 · B2.1 e}}
\lz{115.}{z = $\tfrac{25}{3}$ $\approx$ 8,33}{mit d); Steigungsdreieck 3 : 4 : 5; Abstand 5 ist Kathete, z Hypotenuse: z/5 = $\tfrac{5}{3}$}{}
\end{pfloesung}
\begin{pfloesung}{{\small\color{mbgrau}Abi GK 2024 · B2.1 g}}
\lz{116.}{g(x) = $-$(x $-$ 2) · e$^{-x/2 + 3}$; Durchmesser $\approx$ 7,4 mm}{$f(11) = 9e^{-2{,}5} \approx  0{,}739$; $2 \cdot  0{,}739 LE = 1{,}48 LE = 0{,}74 cm$}{}
\end{pfloesung}
\begin{pfloesung}{}
\lz{117.}{$7$}{$f(3) - g(3) = 10 - 3$}{}
\end{pfloesung}
\begin{pfloesung}{}
\lz{118.}{$40$ cm}{}{}
\end{pfloesung}
\begin{pfloesung}{{\small\color{mbgrau}Abi GK 2026 · A1.4 a}}
\lz{119.}{A = $\tfrac{1}{2}$ · 2$\pi$ · 2 = 2$\pi$}{Grundseite AC = 2$\pi$; Höhe 2}{}
\end{pfloesung}
\begin{pfloesung}{}
\lz{120.}{$5$ m}{Breite $10$ LE $= 20$ m}{}
\end{pfloesung}
\begin{pfloesung}{{\small\color{mbgrau}Abi GK 2023 · B2.1 k}}
\lz{121.}{180 l}{Kopf bei Nr. 23.; mit b); Grundfläche 60 cm $\times$ 60 cm; Höhe 50 cm}{}
\end{pfloesung}
\begin{pfloesung}{{\small\color{mbgrau}Abi GK 2026 · B2.1 f}}
\lz{122.}{passt; Höhe 2 · $\tfrac{4}{3}$ = $\tfrac{8}{3}$ m $\approx$ 2,67 m $\le$ 2,7 m, Breite 2 · (2 + $\tfrac{4}{3}$) = $\tfrac{20}{3}$ m $\approx$ 6,67 m $\le$ 6,7 m}{Halbkreise um ($\pm$2 | 0) mit Radius $\tfrac{4}{3}$ reichen bis x = $\pm\tfrac{10}{3}$}{}
\end{pfloesung}
\begin{pfloesung}{{\small\color{mbgrau}Abi GK 2022 · B2.1 k}}
\lz{123.}{$\approx$ 1041 m}{Kopf bei Nr. 25.; mit j); $h(15) = 759{,}375$; Aufstieg 10 bis 15 min: 900 $-$ 759,4 = 140,6 m}{}
\end{pfloesung}
\begin{pfloesung}{{\small\color{mbgrau}Abi GK 2026 · B2.1 i}}
\lz{124.}{Länge $\approx$ 18,4 m; Zeit $\approx$ 5,25 min $\approx$ 5 min 15 s}{vier Bögen 4 · 2,5 = 10 m; Vollkreis mit Radius $\tfrac{4}{3}$: 2$\pi$ · $\tfrac{4}{3}$ $\approx$ 8,38 m}{}
\end{pfloesung}
\begin{pfloesung}{{\small\color{mbgrau}Abi GK 2024 · B2.1 h}}
\lz{125.}{4,5 cm $\times$ 5,44 cm $\times$ 5,44 cm; Masse $\approx$ 114 g}{mit b); Länge 9 LE = 4,5 cm; Breite und Höhe 2 · 2e $\approx$ 10,87 LE $\approx$ 5,44 cm}{}
\end{pfloesung}
\begin{pfloesung}{}
\lz{126.}{Substitution}{}{}
\end{pfloesung}
\begin{pfloesung}{}
\lz{127.}{$x > 3$}{}{}
\end{pfloesung}
\begin{pfloesung}{}
\lz{128.}{$S_1(-2 | 7)$, $S_2(5 | 14)$}{}{}
\end{pfloesung}
\begin{pfloesung}{}
\lz{129.}{$x > 1$}{}{}
\end{pfloesung}
\begin{pfloesung}{}
\lz{130.}{$x_1 = -3$, $x_2 = 1$}{$2x^2 + 4x - 6 = 0 \Rightarrow x^2 + 2x - 3 = 0$}{}
\end{pfloesung}
\begin{pfloesung}{}
\lz{131.}{$4$ ist die Differenz null: $x < 0$}{$f(x) - g(x) = -x^3 + 8x^2 - 16x = -x \cdot (x - 4)^2$; Grenze $x = 0$, bei $x$}{}
\end{pfloesung}
\begin{pfloesung}{}
\lz{132.}{$t_1 = 2$, $t_2 = 4$}{$h(t) = 4$}{}
\end{pfloesung}
\begin{pfloesung}{}
\lz{133.}{$2$}{}{}
\end{pfloesung}
\begin{pfloesung}{}
\lz{134.}{$x_1 = 2$, $x_2 = -2$}{$f'(x) = -2x + 6$; $-x^2 + 6x - 4 = -2x^2 + 6x$, $x^2 = 4$}{}
\end{pfloesung}
\begin{pfloesung}{{\small\color{mbgrau}Abi GK 2025 · A1.7 a}}
\lz{135.}{f'(x) = $-$2x + 4; x = $-$1 oder x = 1}{$-x^{2} + 4x - 1 = -2x^{2} + 4x \Rightarrow  x^{2} = 1$}{}
\end{pfloesung}
\begin{pfloesung}{{\small\color{mbgrau}Abi GK 2022 · B2.1 d}}
\lz{136.}{(x + 2) · e$^{-x}$ = $-$(x + 1) · e$^{-x}$ $\Rightarrow$ x + 2 = $-$x $-$ 1 $\Rightarrow$ x = $-$1,5}{Kopf bei Nr. 25.}{}
\end{pfloesung}
\begin{pfloesung}{{\small\color{mbgrau}Abi GK 2023 · B2.2 g}}
\lz{137.}{L = (0; 5) $\cup$ (5; $\infty$)}{Kopf bei Nr. 21.; $f(x) - g(x) = 0{,}6x(x - 5)^{2}$}{}
\end{pfloesung}
\begin{pfloesung}{{\small\color{mbgrau}Abi GK 2024 · B2.2 h}}
\lz{138.}{Lösungsweg: |p(x) $-$ f(x)| $\le$ $\tfrac{3}{8}$ ansetzen; für x $\ge$ 1 ist f $\ge$ p $\Rightarrow$ 0,5x$^{4}$ $-$ 0,5x² $\le$ $\tfrac{3}{8}$; Randgleichung mit u = x² lösen (u = 1,5) $\Rightarrow$ 1 $\le$ x $\le$ $\sqrt{1{,}5}$ $\approx$ 1,22}{Kopf bei Nr. 22.; mit g); $0{,}5x^{4} - 0{,}5x^{2} = \tfrac{3}{8} \Rightarrow  x^{2} = 1{,}5$}{}
\end{pfloesung}
\begin{pfloesung}{{\small\color{mbgrau}Abi GK 2025 · B2.2 g}}
\lz{139.}{x $\approx$ 1,9 (grafisch: Stelle, an der der Graph 3 Einheiten weiter rechts 1000 höher liegt); Deutung: Zeitpunkt, zu dem die Anzahl der Likes um 1000 kleiner ist als drei Stunden später}{Kopf bei Nr. 35.}{}
\end{pfloesung}
\begin{pfloesung}{}
\lz{140.}{nach unten}{}{}
\end{pfloesung}
\begin{pfloesung}{}
\lz{141.}{$(0 | 2)$}{}{}
\end{pfloesung}
\begin{pfloesung}{}
\lz{142.}{Punkte eintragen: $(-2 | -6)$, $(-1 | 0)$, $(0 | 0)$, $(1 | 0)$, $(2 | 6)$}{$f(-2) = -6$, $f(-1) = 0$, $f(0) = 0$, $f(1) = 0$, $f(2) = 6$}{}
\end{pfloesung}
\begin{pfloesung}{{\small\color{mbgrau}Abi GK 2025 · B2.2 e}}
\lz{143.}{Anzahl der Likes nimmt zunächst immer stärker zu, ab einem Zeitpunkt immer schwächer, nähert sich 2500}{Kopf bei Nr. 35.}{}
\end{pfloesung}
\begin{pfloesung}{{\small\color{mbgrau}Abi GK 2023 · B2.2 e}}
\lz{144.}{falsch; Wendetangente trifft den Graphen nur in W}{Kopf bei Nr. 21.; mit b), c); Wendetangente y = $-$3x + 14; $f(7) = 20$}{}
\end{pfloesung}
\begin{pfloesung}{}
\lz{145.}{$4{,}5$}{Randwerte $f(-3) = f(3)$}{}
\end{pfloesung}
\begin{pfloesung}{{\small\color{mbgrau}Abi GK 2023 · B2.2 h}}
\lz{146.}{$\approx$ 230 °C (228 bis 235 °C vertretbar)}{Kopf bei Nr. 21.; Schwankung zwischen $\approx$ 222 °C und 245 °C}{}
\end{pfloesung}
\begin{pfloesung}{}
\lz{147.}{$f(3) = 1$ und $f'(3) = 0$}{}{}
\end{pfloesung}
\begin{pfloesung}{}
\lz{148.}{drei Unbekannte $a$, $b$, $c$ und drei Gleichungen}{}{}
\end{pfloesung}
\begin{pfloesung}{}
\lz{149.}{$3x + 2$}{$m = 3$, $n = 2 \Rightarrow y$}{}
\end{pfloesung}
\begin{pfloesung}{}
\lz{150.}{$2x^2 - 4x - 1$}{$c = -1$; $f(1) = a + b - 1 = -3$ und $f'(1) = 2a + b = 0$}{}
\end{pfloesung}
\begin{pfloesung}{}
\lz{151.}{$1$}{$p(2) = g(2)$: $4a + b = 3$; $p'(2) = g'(2)$: $4a = 2$}{}
\end{pfloesung}
\begin{pfloesung}{{\small\color{mbgrau}Abi GK 2022 · B2.2 k}}
\lz{152.}{a = 1/32, b = $\tfrac{1}{8}$; k(0) = 0, k'(0) = 0 (von selbst), k(2) = 1, k'(2) = 1,5}{Kopf bei Nr. 36.; Steigung RS = 1,5; $16a + 4b = 1, 32a + 4b = 1{,}5$}{}
\end{pfloesung}
\begin{pfloesung}{}
\lz{153.}{beide Punkte liegen auf dem Graphen, und durch zwei Punkte geht genau eine Gerade}{$h(0{,}6) = 1{,}2$ und $h(3) = 6$}{}
\end{pfloesung}
\begin{pfloesung}{{\small\color{mbgrau}Abi GK 2022 · B2.2 g}}
\lz{154.}{Steigung (1,2 + 0,4)/(2 $-$ 0,4) = 1, y-Achsenabschnitt $-$0,8 $\Rightarrow$ g(x) = x $-$ 0,8; beide Punkte erfüllen die Gleichung}{Kopf bei Nr. 36.}{}
\end{pfloesung}
\begin{pfloesung}{{\small\color{mbgrau}Abi GK 2023 · B2.2 k}}
\lz{155.}{k(x) = 2x³ $-$ 28x² + 130x + 20}{Kopf bei Nr. 21.; $d = 20, c = 130$; $125a + 25b = -450, 75a + 10b = -130$}{}
\end{pfloesung}
\begin{pfloesung}{}
\lz{156.}{die Randwerte $A(0)$ und $A(4)$}{}{}
\end{pfloesung}
\begin{pfloesung}{}
\lz{157.}{$60$}{}{}
\end{pfloesung}
\begin{pfloesung}{}
\lz{158.}{$6u - u^3$}{$A(u) = u \cdot f(u) = u \cdot (6 - u^2)$}{}
\end{pfloesung}
\begin{pfloesung}{}
\lz{159.}{$x = 2$}{$A'(x) = 12 - 3x^2$; $12 - 3x^2 = 0 \Rightarrow x^2 = 4$}{}
\end{pfloesung}
\begin{pfloesung}{}
\lz{160.}{$28 - 2a$}{$2a + b = 28$}{}
\end{pfloesung}
\begin{pfloesung}{}
\lz{161.}{$\tfrac{1}{4}$}{$d(x) = p(x) - f(x) = x^2 - x^4$}{}
\end{pfloesung}
\begin{pfloesung}{}
\lz{162.}{$P(3 | 1{,}34)$}{$A(u) = u \cdot f(u) = (u^2 + 3u) \cdot e^{-0{,}5u}$}{}
\end{pfloesung}
\begin{pfloesung}{{\small\color{mbgrau}Abi GK 2024 · B2.2 g}}
\lz{163.}{d(x) = 0,5x² $-$ 0,5x$^{4}$ maximal bei x = $\pm$1/$\sqrt{2}$ mit d = $\tfrac{1}{8}$}{Kopf bei Nr. 22.; $d'(x) = x - 2x^{3} = 0 \Rightarrow  x = 0, \pm 1/\sqrt{2}$; $d(0) = d(\pm 1) = 0$}{}
\end{pfloesung}
\begin{pfloesung}{{\small\color{mbgrau}Abi GK 2022 · B2.1 h}}
\lz{164.}{P($\sqrt{2}$; ($\sqrt{2}$ + 2) · e$^{-√2}$) $\approx$ (1,41; 0,83)}{Kopf bei Nr. 25.; $A(u) = u \cdot  f(u) = u(u + 2)e^{-u}$; $A'(u) = (2 - u^{2})e^{-u} = 0 \Rightarrow  u = \sqrt{2}$}{}
\end{pfloesung}
\end{document}